\section{Introduction}
\label{sec:introduction}

Executable RTL benchmarks have made functional correctness measurable for
language-model-generated hardware~\cite{liu2023verilogeval,lu2024rtllm}.
Physical usefulness is a different capability.  Behaviorally compatible modules can
implement the same computation with sharply different critical paths,
resource mappings, and achievable clocks.  A model trained only to pass a
testbench has no reason to prefer the physically better form.  Conversely,
generating many candidates and selecting the fastest after electronic-design-
automation (EDA) evaluation shifts the cost to inference and repeated physical
design.

Recent work has therefore connected RTL generation to synthesis feedback,
reinforcement learning, evolutionary search, or per-design adaptation
~\cite{ppartl2025,chipseek2025,yao2024rtlrewriter,poet2026,
evolve2026,zhou2026alphartl,drrtl2026}.  A central measurement problem remains:
an optimizer can improve a selected candidate, a conditional-on-correct mean,
or the learned objective without making a random policy draw more useful.
Physical implementation is also too expensive for every policy update, which
makes a learned reward attractive but exposes it to distribution shift and
overoptimization~\cite{gao2023scaling}.  The reward is least trustworthy where
the optimized policy places its mass.

The open question in this paper is therefore deliberately bounded: within a
controlled streaming-accelerator domain, does proxy-guided training change the
generator's physical utility at a fixed draw budget after functional and
implementation failures are charged to the policy that produced them?
Answering that question also requires distinguishing a timing-aware reward from
generic policy movement and from an increase in correctness alone.  The group-
relative update rule and EDA feedback are ingredients of the study rather than
new algorithms; the contribution is the failure-aware distribution-level
formulation, controlled reward-package comparisons, and a prospective post-route test.

We address these two issues with a strict training/evidence boundary
(Fig.~\ref{fig:training-boundary}).  During training, an executable oracle
first rejects functionally invalid RTL; a learned physical reward ranks only
admissible completions.  During reporting, neither the reward nor a selected
example is accepted as hardware evidence.  Every distinct correct candidate is
implemented by the same post-route Vivado flow, its original sample
multiplicity is restored, and every incorrect or implementation-failed draw is
assigned zero frequency.  The resulting equal-sample endpoint measures the
physical utility of one draw while keeping the design, rather than the
candidate, as the unit of aggregation.

\begin{figure*}[!t]
\centering
\includegraphics[width=\textwidth]{fig01_training_evidence_boundary_v3}
\caption{Implemented training and physical-evaluation workflow.
(a) Domain SFT and the physical predictor use training-only resources.
(b) Fresh policy samples receive group-relative updates with a frozen SFT
reference; flat groups do not advance the update counter.
(c) The RF reward uses canonicalized structural features; MLP and
correctness-only rewards are separate control runs. Red branches assign zero
to rejected candidates.
(d) Frozen checkpoints are evaluated using the executable oracle and post-route
implementation. Original multiplicity and the full draw denominator are
restored; incorrect and implementation-failed draws are not discarded.
Dashed paths denote feedback, references, or separately run controls.}
\label{fig:training-boundary}
\end{figure*}

The primary test is a separately frozen
\studyclaim{StudyTwoDesignCount}-design study spanning interpolation and
extrapolation across the retained accelerator generators.  Two policy-training
seeds sharing one SFT checkpoint and a canonicalization-aware random-forest (RF) reward raise
penalized equal-sample \fmax{} from
\studyclaim{StudyTwoOverallSftPenalized} to
\studyclaim{StudyTwoOverallRfPenalized}~MHz.  Their seed-averaged paired effect
is \studyclaim{StudyTwoOverallRfSftGain}~MHz with interval
[\studyclaim{StudyTwoOverallRfSftCiLow},
\studyclaim{StudyTwoOverallRfSftCiHigh}]~MHz; each training seed is positive,
and the interpolation and extrapolation intervals both exclude zero.  The
optimized policy's correctness remains close to the supervised fine-tuned
(SFT) baseline, and the failure penalty already charges the endpoint for the
difference.

Matched controls constrain the explanation for the gain.  The repaired RF reward
exceeds the original multilayer-perceptron (MLP) reward by
\studyclaim{StudyTwoOverallRfMlpGain}~MHz under the same tasks and sampling
budget.  Correctness-only training instead raises validity to
\studyclaim{StudyTwoCorrectnessOnlyCorrectness}\% but lowers penalized \fmax{}
relative to SFT in its single preregistered seed.  The observed physical gain
therefore does not arise from higher functional yield; attribution beyond these
realized checkpoints requires further training replication.  Conditional resource and power
analysis exposes its implementation cost.  A separate, explicitly descriptive
PYNQ-Z2 sweep is retained in the evidence record rather than used to strengthen
the primary population claim.

This paper makes four contributions:

\begin{itemize}
    \item \emph{Controlled empirical finding:} on the frozen accelerator split,
    both independently trained RF policies improve the failure-penalized
    post-route endpoint under interpolation and extrapolation.  Descriptive
    family and sensitivity analyses expose, rather than average away, the weak
    family.
    \item \emph{Reward-package comparison:} original-reward and correctness-only
    arms separate the observed gain from an increase in functional yield.
    The comparison supports the repaired reward package on the studied split;
    it does not isolate random-forest architecture from canonical features.
    \item \emph{Failure-aware protocol:} the fixed-draw estimand preserves
    original sample multiplicity and charges every incorrect or implementation-
    failed draw zero, while real post-route EDA rather than the learned proxy
    remains the reporting authority.
    \item \emph{Bounded engineering interpretation:} candidate probabilities,
    a source-level RTL contrast, conditional PPA, setup-closure and board
    sensitivities, and matched general-RTL evaluation expose both physical
    improvements and specialization costs.
\end{itemize}

The dominant observed mechanism is probability reallocation, not demonstrated
open-ended circuit invention.  Supervision exposes multiple correct
implementations, and policy optimization often makes fast forms more probable
when the reward can distinguish them.  The sealed support diagnostic also has
exceptions, so the paper does not reduce every gain to reweighting within the
finite SFT sample.  This narrower interpretation predicts stronger gains where
fast structures are representable and reward-distinguishable, weaker families,
and failure when a learned proxy ceases to track physical reality.
